\begin{frame}{자주 하는 실수: $\gamma$ 를 ``느낌''으로 고르기}
  \begin{itemize}
    \item \texttt{gamma}(대략 $\frac{1}{2\sigma^2}$)가 ``결정
      경계가 얼마나 구불구불해질 수 있는가''를 \textbf{통째로}
      결정.
    \item \kb{$\gamma$ 극단적으로 크게}($\sigma$ 작음): 각 학습점
      바로 주변에 좁고 뾰족한 ``영향권'' $\Rightarrow$ 학습 데이터
      하나하나를 거의 다 따로 감싸는 \textbf{심하게 구불거림}
      (전형적 \textbf{과적합}).
    \item \kb{$\gamma$ 너무 작게}: 모든 점의 영향권이 서로 겹쳐
      사실상 \textbf{선형 커널과 비슷}하게 뭉뚱그려진 경계만
      나옴(과소적합).
  \end{itemize}
  \vskip0.4em
  \begin{block}{수치적 증거: 앞 $\gamma$ 스윕 표}
    $\gamma=0.05$: test 0.467로 \textbf{동전 던지기보다 나쁨}(선형
    0.517보다도 낮음) — RBF가 ``거의 상수 행렬''에 가까워져
    ``모든 점이 비슷'' $\Rightarrow$ 사실상 \textbf{무작위}.
    $\gamma=30$: train/test 1.00 — \textbf{이 데이터는 노이즈가
    적어} 운 좋게 일반화되는 것. 노이즈 있는 실전 데이터에선
    개별 노이즈 패턴까지 따라붙어 test가 \textbf{급락}.
  \end{block}
  \vskip0.4em
  {\footnotesize 4.1절 kNN의 $k$, 5.2절의 $C$와 \textbf{같은 종류}의
  ``복잡도 손잡이'' — $C$와 $\gamma$ 를 \textbf{동시에} 그리드
  서치하는 것이 표준(서로 영향을 주고받음).}
  \vskip0.3em
  \begin{center}
  \kq{검증셋으로 $(C, \gamma)$ 를 함께 그리드 서치하는 것 =\\
      5.2절 $C$ 튜닝의 비선형 버전.}
  \end{center}
\end{frame}
